% !TeX root = ../main.tex
% Fonte: Lezione 14, pp. 7--12. Precisate le ipotesi di dominio e completato l'ultimo esercizio.
\section{Due limiti trigonometrici fondamentali}
\label{sec:limiti-trigonometrici}

La continuità permette di calcolare $\sin x_n\to0$ e $\cos x_n\to1$ quando $x_n\to0$. Non risolve però direttamente i rapporti in cui numeratore e denominatore tendono entrambi a zero. In questi casi servono informazioni più precise sulla loro grandezza relativa.

\begin{teorema}
Se $x_n\to0$ e $x_n\ne0$ definitivamente, con gli angoli misurati in radianti, allora:
\begin{enumerate}[label=\arabic*.]
\item $\sin x_n/x_n\to1$.
\item $(1-\cos x_n)/x_n^2\to1/2$.
\end{enumerate}
\end{teorema}
\begin{proof}
\begin{enumerate}[label=\arabic*.]
\item \textbf{Il rapporto con il seno.} Per $0<t<\pi/2$, il confronto geometrico fra un triangolo inscritto, un settore circolare e un triangolo tangente dà
\[
\sin t<t<\tan t.
\]
Più precisamente, nel cerchio unitario le tre aree sono $\sin t/2$, $t/2$ e $\tan t/2$. Dividendo per $\sin t>0$ si ricava
\[
1<\frac{t}{\sin t}<\frac1{\cos t},
\qquad
\cos t<\frac{\sin t}{t}<1.
\]
Il rapporto $\sin t/t$ è pari, quindi la stessa stima vale per $0<|t|<\pi/2$. Definitivamente la si può applicare a $t=x_n$. La continuità del coseno dà $\cos x_n\to1$; per i carabinieri, $\sin x_n/x_n\to1$.
\item \textbf{Il rapporto con il coseno.} Poiché $\cos x_n\to1$, definitivamente $1+\cos x_n\ne0$. Moltiplicando per il coniugato e usando $1-\cos^2t=\sin^2t$,
\[
\frac{1-\cos x_n}{x_n^2}
=\frac{\sin^2x_n}{x_n^2(1+\cos x_n)}
=\left(\frac{\sin x_n}{x_n}\right)^2\frac1{1+\cos x_n}
\longrightarrow1\cdot\frac12.
\]
\end{enumerate}
\end{proof}

\begin{schema}
\begin{tikzpicture}[scale=2.8]
  \draw[asse] (-1.2,0) -- (1.5,0);
  \draw[asse] (0,-1.12) -- (0,1.25);
  \draw[black!55] (0,0) circle[radius=1];
  \draw[intervallo] (0,0) -- (1,{tan(35)});
  \draw[black!55] (0,0) -- (1,0);
  \draw[intervallo] (1,-.15) -- (1,.95);
  \draw[black!55] (1,0) -- ({cos(35)},{sin(35)});
  \draw[dashed,coloreschema] ({cos(35)},0) -- ({cos(35)},{sin(35)});
  \draw[coloreschema] (.3,0) arc[start angle=0,end angle=35,radius=.3];
  \node at (.4,.12) {$t$};
  \node[below left] at (0,0) {$O$};
  \node[below right] at (1,0) {$A$};
  \node[above left] at ({cos(35)},{sin(35)}) {$B$};
  \node[above right] at (1,{tan(35)}) {$T$};
\end{tikzpicture}
\par\smallskip Nel cerchio unitario, $B=(\cos t,\sin t)$ e $T=(1,\tan t)$. Il triangolo $OAB$ è contenuto nel settore $OAB$, a sua volta contenuto nel triangolo $OAT$.
\end{schema}

\begin{esercizio}[Riconoscere gli argomenti dei limiti notevoli]
Calcolare
\[
\lim_{n\to+\infty}\frac{\sin^2(3/n)-1/n^2}{\cos(2/n)-1}.
\]
\end{esercizio}
\begin{svolgimento}
Numeratore e denominatore tendono a zero. Per $n\ge1$ il denominatore è non nullo, perché $0<2/n\le2<2\pi$ e $\cos(2/n)<1$. Riscriviamo separatamente
\[
\sin^2(3/n)-\frac1{n^2}
=\frac1{n^2}\left[9\left(\frac{\sin(3/n)}{3/n}\right)^2-1\right]
\]
e
\[
\cos(2/n)-1
=-\frac4{n^2}\,\frac{1-\cos(2/n)}{(2/n)^2}.
\]
Il fattore $1/n^2$ si semplifica esattamente, prima di passare al limite. Si ottiene
\[
\frac{9\left(\frac{\sin(3/n)}{3/n}\right)^2-1}
{-4\,\frac{1-\cos(2/n)}{(2/n)^2}}
\longrightarrow\frac{9-1}{-4\cdot\frac12}=-4.
\]
\end{svolgimento}

\section{Equivalenza e cancellazioni}
% Fonte: Lezione 14, pp. 11--12. L'esito finale dell'esempio non è calcolato nella fonte.

\begin{definizione}[Successioni asintoticamente equivalenti]
Se $b_n\ne0$ definitivamente, si scrive $a_n\sim b_n$ per indicare che $a_n/b_n\to1$. In particolare, se $x_n\to0$ e $x_n\ne0$ definitivamente, $\sin x_n\sim x_n$.
\end{definizione}

La relazione $\sin x_n\sim x_n$ riguarda un rapporto: non è un'uguaglianza. Moltiplicare e dividere per l'argomento consente di usare il limite notevole senza cambiare l'espressione. Sostituire direttamente due termini che poi si sottraggono può invece cancellare proprio l'informazione decisiva.

\begin{esercizio}[Una cancellazione fra infinitesime]
Calcolare
\[
\lim_{n\to+\infty}\frac{\sin(1/n)-1/n}{1/n^4},
\]
spiegando perché non si possa sostituire semplicemente $\sin(1/n)$ con $1/n$ nel numeratore.
\end{esercizio}
\begin{svolgimento}
Il raccoglimento esatto dà
\[
\frac{\sin(1/n)-1/n}{1/n^4}
=n^3\left[\frac{\sin(1/n)}{1/n}-1\right].
\]
La parentesi tende a zero, ma il fattore $n^3$ diverge: rimane una forma indeterminata. La sola equivalenza $\sin(1/n)\sim1/n$ non autorizza a cancellare i due addendi e a concludere che il limite sia zero.

\textbf{Una stima più precisa.} Poniamo $t=1/(3n)$. Dall'identità $\sin(3t)=3\sin t-4\sin^3t$, conseguenza delle formule di addizione, segue
\[
\frac1n-\sin\frac1n
=3(t-\sin t)+4\sin^3t\ge4\sin^3t,
\]
perché $0<\sin t<t$. Quindi
\[
\frac{\sin(1/n)-1/n}{1/n^4}
\le-4n^4\sin^3\!\left(\frac1{3n}\right)
=-\frac{4n}{27}\left(\frac{\sin(1/(3n))}{1/(3n)}\right)^3
\longrightarrow-\infty.
\]
Il confronto con una successione tendente a $-\infty$ conclude che il limite cercato è $-\infty$.
\end{svolgimento}
